\section{Apéndice: ruta de aprendizaje de Lean y AI for Math}

\noindent\begin{longtable}{p{0.20\textwidth}p{0.35\textwidth}p{0.36\textwidth}}
\toprule
Etapa & Qué aprender & Por qué aprenderlo \\
\midrule
Lenguaje matemático & Proposiciones, definiciones, lemma, proof, contraejemplos & Entender cómo se organiza una demostración en lenguaje natural. \\
Fundamentos de Lean & theorem, example, rw, simp, apply, exact, induction & Entender cómo se interactúa con un proof assistant. \\
Mathlib & Biblioteca de teoremas de uso común, nomenclatura, búsqueda, relaciones de dependencia & Las demostraciones con IA dependen de las bibliotecas matemáticas existentes. \\
Autoformalization & Del lenguaje natural a un statement de Lean & Conectar artículos y enunciados de problemas con el sistema formal. \\
Proof Search & Combinación de tactics, goal state, corrección de errores & Permitir que el modelo itere con la retroalimentación del verificador. \\
Human Collaboration & Criterio estético de los matemáticos, elección de problemas, capacidad explicativa & Juzgar si una demostración tiene sentido, y no solo si pasa la verificación. \\
\bottomrule
\end{longtable}

\begin{knowledgebox}{El sentido de la ruta de aprendizaje}
Esta tabla cambia el orden de aprendizaje de AI for Math: en lugar de «primero aprender los modelos», propone «primero entender las matemáticas y los sistemas formales». Para los lectores con formación en ingeniería, los fundamentos de Lean y el proof search son el requisito mínimo para entrar en este campo; para los lectores con formación matemática, la autoformalization y la retroalimentación de los modelos son la vía de entrada para entender los límites de las capacidades de la IA.
\end{knowledgebox}

\subsection{Resumen de la sección}

AI for Math es una ruta interdisciplinaria: el lenguaje matemático, los sistemas formales, el entrenamiento de modelos y la colaboración entre humanos y máquinas son todos indispensables. Cuanto más clara sea la ruta de aprendizaje, menos fácil será malinterpretarla como «hacer que el modelo resuelva cada vez más ejercicios».
